\documentclass[a4paper, 10pt]{article}
\usepackage[margin=2cm]{geometry}
\usepackage{graphicx}
\usepackage{amsmath}
\usepackage{booktabs}
\usepackage{parskip}
\usepackage{enumitem}
\renewcommand{\familydefault}{\ttdefault}
\usepackage[scale=0.9]{FiraMono}
\usepackage{url}

\title{\textbf{CSU11022 Mid-Semester Test Cheat Sheet}}

\begin{document}
\maketitle

% ─────────────────────────────────────────────
\section{Two's Complement for Signed Numbers}

\begin{itemize}
  \item For 8 bits, the range is $-128$ to $+127$.
  \item Any value with a 1 in the leftmost bit is negative; any value with a 0 is positive.
  \item Conversion: flip every bit, then add 1 (converts between positive and negative signed values).
\end{itemize}

\textbf{Example: What is $-4$ in 8-bit two's complement?}

\begin{enumerate}
  \item Find 4 in binary:\quad \texttt{0000 0100}
  \item Flip all bits:\quad \texttt{1111 1011}
  \item Add 1:
\begin{verbatim}
  1111 1011
+ 0000 0001
-----------
  1111 1100
\end{verbatim}
\end{enumerate}

% ─────────────────────────────────────────────
\section{Condition Code Flags}

These flags are only modified by a \texttt{CMP} operation or an arithmetic operation
with an \texttt{s} suffix. \texttt{CMP} acts as subtraction but does not write to a
register --- it only updates the flags.

\begin{itemize}
  \item \textbf{C (Carry):} Set to 1 if the operation produces a carry out of the most
        significant bit (e.g.\ adding two 8-bit numbers where the result would require
        a 9th bit); otherwise 0.
  \item \textbf{Z (Zero):} Set to 1 if all bits of the result are zero; otherwise 0.
  \item \textbf{N (Negative):} Set to 1 if the most significant bit of the result is 1;
        otherwise 0.
  \item \textbf{V (oVerflow):} Set to 1 if both operands have the same sign but the
        result has a different sign; otherwise 0.
\end{itemize}

% ─────────────────────────────────────────────
\section{Arrays}

\begin{itemize}
  \item Row-major ordering (row index comes first).
  \item To find the address offset for an element:
        \begin{align*}
          \text{index}        &= (\text{row} \times \text{row\_size}) + \text{col} \\
          \text{byte\_offset} &= \text{index} \times \text{elem\_size}             \\
          \text{address}      &= \text{array\_base\_address} + \text{byte\_offset}
        \end{align*}
\end{itemize}

% ─────────────────────────────────────────────
\section{How the Stack Works}
\begin{itemize}
  \item \textbf{PUSH} --- Places items on top of the stack.
        Equivalent to: \texttt{STR R0, [R13, \#-4]!}
  \item \textbf{POP} --- Removes item from top of stack.
        Equivalent to: \texttt{LDR R0, [R13], \#4}
  \item The stack grows downwards in ARM assembly (\emph{full descending}).
  \item LIFO (Last In, First Out).
  \item The highest-numbered register in curly braces is always pushed/popped first.
  \item Everything pushed onto the stack must be popped off.
\end{itemize}

% ─────────────────────────────────────────────
\section{Subroutines}
\begin{itemize}
  \item Parameters are passed using registers \texttt{R0}--\texttt{R3}.
        \begin{itemize}
          \item If more parameters are needed, use the system stack:
                \begin{itemize}
                  \item Calling program pushes parameters onto the stack.
                  \item Subroutine accesses parameters relative to the stack pointer.
                  \item Calling program pops parameters off the stack after the
                        subroutine returns.
                \end{itemize}
          \item \texttt{R0}--\texttt{R3} should be assumed clobbered by the subroutine
                unless used to pass data in/out.
        \end{itemize}
  \item If calling another subroutine from within a subroutine, \textbf{PUSH} the link
        register \texttt{LR} at the start.
  \item Subroutines must \textbf{PUSH} all registers they use onto the system stack at
        the \emph{start} of the subroutine, and \textbf{POP} them back at the \emph{end}.
  \item Subroutines can call themselves recursively (e.g.\ computing powers).
  \item All labels except \texttt{main} must be prefixed with \texttt{.L}.
\end{itemize}

% ─────────────────────────────────────────────
\section{Handy Reference}

\subsection*{Hexadecimal Table}

\begin{center}
\begin{tabular}{ccc}
  \toprule
  Hex & Decimal & Binary \\
  \midrule
  0   & 0       & 0000   \\
  1   & 1       & 0001   \\
  2   & 2       & 0010   \\
  3   & 3       & 0011   \\
  4   & 4       & 0100   \\
  5   & 5       & 0101   \\
  6   & 6       & 0110   \\
  7   & 7       & 0111   \\
  8   & 8       & 1000   \\
  9   & 9       & 1001   \\
  A   & 10      & 1010   \\
  B   & 11      & 1011   \\
  C   & 12      & 1100   \\
  D   & 13      & 1101   \\
  E   & 14      & 1110   \\
  F   & 15      & 1111   \\
  \bottomrule
\end{tabular}
\end{center}

\subsection*{MOV vs LDR}

\texttt{MOV} can only store a limited range of values, whereas \texttt{LDR} does not
have this limitation.

% ─────────────────────────────────────────────
\section{Definitions}

\begin{itemize}
  \item \textbf{bit} --- 1 binary digit
  \item \textbf{byte} --- 8 binary digits
  \item \textbf{halfword} --- 2 bytes
  \item \textbf{word} --- 4 bytes
  \item \textbf{LSL} --- Logical Shift Left (equivalent to multiplying by $2^{n}$)
  \item \textbf{LSR} --- Logical Shift Right (equivalent to dividing by $2^{n}$)
  \item \textbf{LDM / STM} --- Load/Store Multiple. Requires word alignment
        (\texttt{memoryAddress \% 4 == 0})
  \item \textbf{LDMIA / STMIA} --- IA = Increment After. Automatically increments the
        base address by the number of bytes read/written.
\end{itemize}

% ─────────────────────────────────────────────
\section{Efficiency Notes}

\begin{itemize}
  \item Every instruction takes at least one cycle.
  \item \texttt{UDIV} requires multiple cycles (dependent on the values involved).
  \item Each memory access costs one extra cycle.
        Relevant for: \texttt{LDR}, \texttt{STR}, \texttt{LDM}, \texttt{STM},
        \texttt{PUSH}, \texttt{POP}.
  \item Reading words or halfwords across a word boundary costs \emph{two} extra cycles
        instead of one.
  \item Taken branches incur a 2-cycle pipeline stall penalty.
        \begin{itemize}
          \item Ensure you actually need the branch.
          \item If-then-else constructs are expensive.
          \item For conditional branches, optimise for the common case.
        \end{itemize}
\end{itemize}

\bigskip

\begin{center}
\bigskip
\Large\bfseries License
\end{center}

\bigskip

Copyright \copyright\ 2026 Tommy Byrne

\bigskip

This work is free documentation: you can redistribute it and/or modify
it under the terms of the GNU General Public License as published by
the Free Software Foundation, either version 3 of the License, or
(at your option) any later version.

\bigskip

This work is distributed in the hope that it will be useful,
but WITHOUT ANY WARRANTY; without even the implied warranty of
MERCHANTABILITY or FITNESS FOR A PARTICULAR PURPOSE. See the
GNU General Public License for more details.

\bigskip

You should have received a copy of the GNU General Public License
along with this work. If not, see
\url{https://www.gnu.org/licenses/}.

\bigskip

Source files are available at:
\url{https://github.com/mbdalpha/openNotes}

\end{document}